\chapter{Optic Neuritis}
\index{Optic Neuritis}

\medskip
\noindent\textit{\textbf{Under review.} This chapter is an AI-assisted educational draft; it is not a substitute for professional medical advice.}

\section*{Definition and Overview}
Optic neuritis is inflammation or immune-mediated injury of the optic nerve, causing visual symptoms. It may occur alone or with conditions such as multiple sclerosis, neuromyelitis optica spectrum disorder, or MOG-antibody disease.

\section*{Clinical Features}
Typical features are subacute loss of vision in one eye, reduced colour vision, a central blind spot, and pain worsened by eye movement. Symptoms may affect both eyes or recur, particularly in atypical cases.

\section*{When to Seek Medical Care}
New visual loss, severe eye pain, or loss of colour vision needs urgent medical assessment. Contact your local emergency services for sudden neurological symptoms such as facial droop, limb weakness, speech difficulty, or collapse.

\section*{Diagnosis and Investigations}
Assessment includes visual acuity, colour vision, pupils, eye movements, fundus examination, visual fields, and optical coherence tomography. MRI of the brain and orbits and blood tests for alternative inflammatory or infectious causes may be needed.

\section*{Management}
Treatment depends on severity and cause. Specialist care may include intravenous or oral corticosteroids for selected attacks and disease-specific treatment or long-term monitoring when an associated inflammatory or demyelinating disorder is found.

\section*{Prognosis and Self-Care}
Vision often improves over weeks to months, but residual loss or recurrence can occur. Attend follow-up, report new neurological or visual symptoms promptly, and do not drive if vision is unsafe.
\section*{Safety-netting}
Seek urgent care for sudden severe symptoms, chest pain, breathing difficulty, collapse, new neurological deficit, or rapidly worsening function.
